\section{Tagged paths and compatible barriers}\label{sec:barriers}

Fix a countable ordered list of coordinates satisfying \textup{(P)}, with
flag set $\Sigma$.  Throughout this section the list is fixed.  We seek a
bounded smooth real coordinate that vanishes on a compact testing set
in $G$ and equals one on a collar about a chosen dust point.  The
coordinate must remain compatible with the prior list.  A lower bound on
path costs will supply the separation, while the norms of the density and
boundary tolls will control the barrier's size.

\subsection{Endpoint data and barriers}

The endpoint data specify the two sets to be separated.  Gates optionally
restrict where paths cross the prescribed collars; the controlled energy
then omits a permitted penalty region.  In the ungated case this region
is empty.

\begin{definition}[Endpoint data and a barrier]\label{def:endpoint}
An endpoint problem consists of a point $b\in B$, a nested sequence
of prescribed dust disks $U_j=U_j(b)$ with boundaries $C_j$, and a
compact anchor $E\subset G$ with nonempty interior, disjoint from
all $\overline U_j$.  A compact set $E_*\subset\operatorname{int}E$
is also specified.  The collar sequence may start at any depth in the
fixed exhaustion.

There may additionally be \emph{gates}: a fixed normalized univalent
map $f_0$ on $G$, a bounded open Euclidean ball $W_2$, and a closed
ball $W_1\subset W_2$.  Every hit of every $C_j$ must then lie in
$f_0^{-1}(\overline W_2)$.  In this case a \emph{penalty region} is
any relatively compact open $O\subset G$ such that
\[
 \overline O\cap E=\varnothing,
 \qquad f_0(\overline O)\cap W_1=\varnothing.
\]
In the absence of gates put $O=\varnothing$.

A barrier is a bounded smooth real coordinate $u$ on $G$, with finite
total energy and budget $a_u\in\ell^2(B)$, such that $u=0$ on $E_*$,
$u=1$ on some $C_l$, and appending $u$ preserves \textup{(P)} after a
countable enlargement of the flag set.  Its controlled size is
\[
 \|\dd u\|_{2,G\setminus O}+\|a_u\|_{\ell^2(B)}.
\]
Energy on $O$ is not included in this size, but must still be finite.
\end{definition}

\subsection{Path length, tags, and residual contacts}

Every path has parameter interval $[0,1]$ and is continuous as a map
$\gamma:[0,1]\to X$.  On the open set $\gamma^{-1}(G)$ it is a path
in the physical sphere.  Its \emph{domain-part length} is the sum of
its physical spherical lengths on the components of this open set,
where the length on an open component is obtained by exhaustion by
compact subintervals.  Only paths with finite domain-part length will
be used.  Their arclength occupation measure on $G$ is denoted by
$\nu_\gamma$; it counts length with parameter multiplicity.
When an integral is written over parameter times, $\dd s$ means the
path's domain-part arclength measure on those times, not Lebesgue
measure $\dd t$; its pushforward to $G$ is $\nu_\gamma$.

We call a pair $(g,m)$ a \emph{metric} when $g\in L^2(G,\dd a)$
is nonnegative Borel, possibly with infinite values on a null set,
and $m\in\ell^2(B)$ is nonnegative.  The use of a domain
density together with square-summable boundary weights is in the
tradition of Schramm's transboundary extremal length
\cite{Schramm1995}.  Here a boundary toll is charged at every visit
time.  This convention will be used in the concatenation and
intersection arguments.  The true-path cost is
\begin{equation}\label{eq:true-path-cost}
 \mathcal L_{g,m}(\gamma)
 =\int_G g\,\dd\nu_\gamma
   +\sum_{r\in\gamma^{-1}(B)}m(\gamma(r)).
\end{equation}
The second term counts \emph{times}, not merely distinct image points.
As elsewhere, its value is the supremum of finite subsums.  Thus a
positive toll at a point visited infinitely many times gives infinite
cost.  We write $\mathcal L^-_{g,m}$ when the time $1$ is omitted from
the toll sum, and $\mathcal L^\circ_{g,m}$ when both times $0$ and $1$
are omitted.  These omissions do not affect the arclength term.

Arclength on the open domain times is atomless.  Its occupation
dominates $\HH^1_{\mathrm{sph}}$ on the image in $G$.  If $\gamma$ is
injective and its image in $G$ has finite physical $\HH^1$ measure,
then its domain-part length is finite and its occupation is precisely
that image measure.  To check these statements it suffices to exhaust
the open domain times by compact intervals.  On each such interval
one has the ordinary facts for continuous rectifiable curves.  For an
injective curve the partition sums are bounded by its image length,
because the intervening subarcs have disjoint interiors.  In particular,
all uses of curve coarea below concern ordinary rectifiable portions
in the physical domain, irrespective of the behavior at dust times.

We next record portions of a path that travel inside finite unions of
prior level sets.  A tag will specify such a union and the parameter
interval on which the path follows it.  The dust contacts outside
these intervals will be subject to separate restrictions.

A \emph{type} is a nonempty finite list
$\boldsymbol\iota=(i_1,\ldots,i_d)$ of prior labels, with repetitions
allowed.  Write
\[
 H_{\boldsymbol\iota}(\mathbf t)
   =\bigcup_{j=1}^d H_{i_j}(t_j),
 \qquad \mathbf t=(t_1,\ldots,t_d)\in\T^d.
\]
There are countably many types.  For each type fix a Haar-null Borel
set $N_{\boldsymbol\iota}\subset\T^d$ of excluded tuples; it may be
empty.  A path has countably many optional \emph{tag slots}.  An active
slot specifies a type, a tuple outside its excluded set, and an interval
$(a,b)$ with $0\le a<b\le1$, such that $\gamma$ is nonconstant on
$(a,b)$ and its image there is contained in
$H_{\boldsymbol\iota}(\mathbf t)$.  The entire interval $(0,1)$ is
permitted, and different tags may overlap.  For a tagged path
$\Gamma=(\gamma,\theta)$, define the compact set of residual dust times
by
\[
 Z_\Gamma
   =\gamma^{-1}(B)\setminus
       \bigcup_{\text{active slots}}(a,b).
\]

The additional data defining a path system are as follows:
\begin{itemize}
\item a finite atomless Borel measure $\lambda$ on $B$, possibly zero;
\item a countable collection $\mathscr Q$ of earlier probability laws
on paths, for each of which
\[
 m_\xi(s)=\Prob_\xi\{s\in\gamma([0,1])\},\qquad s\in B,
 \quad\text{satisfies}\quad m_\xi\in\ell^2(B);
\]
\item increasing finite sets $P_{\xi,n}$ exhausting
$\{m_\xi>0\}$, and increasing finite sets $P_{\Sigma,n}$ exhausting
$\Sigma$.
\end{itemize}
Finite exhaustions may stabilize.  Laws carrying their own tags are
evaluated on their path marginals in the displayed definition.

The residual image must have zero $\lambda$ measure.  Its incidences
with earlier path laws and prior levels will be bounded by the
corresponding point weights, summed over residual times, on each
rational test interval and after each prescribed finite deletion.

\begin{definition}[An allowed tagged path]\label{def:path-system}
A finite-domain-length tagged path is \emph{allowed} if all of its tags
obey their tuple exclusions, if
\begin{equation}\label{eq:path-atomless-constraint}
 \lambda\bigl(\gamma(Z_\Gamma)\bigr)=0,
\end{equation}
and if the following inequalities hold for every closed interval
$I\subset[0,1]$ with rational endpoints, including $I=[0,1]$:
\begin{align}
 &\Prob_\xi\{\text{the sampled path visits }
           \gamma(Z_\Gamma\cap I)\setminus P_{\xi,n}\}
 \nonumber\\[-2pt]
 &\hspace{35mm}\le
   \sum_{\substack{r\in Z_\Gamma\cap I\\
                   \gamma(r)\notin P_{\xi,n}}}
              m_\xi(\gamma(r))
 \qquad(\xi\in\mathscr Q,\ n\in\N),
 \label{eq:path-law-constraint}\\
 &\left|\left\{v\in\T:
      H_i(v)\cap
       \bigl(\gamma(Z_\Gamma\cap I)\setminus P_{\Sigma,n}\bigr)
                      \ne\varnothing\right\}\right|
 \nonumber\\[-2pt]
 &\hspace{35mm}\le
   \sum_{\substack{r\in Z_\Gamma\cap I\\
                   \gamma(r)\notin P_{\Sigma,n}}}
              a_i(\gamma(r))
 \qquad(i\text{ a prior label},\ n\in\N).
 \label{eq:path-level-constraint}
\end{align}
An untagged path is called allowed if it admits such tags.  Unless
otherwise indicated, path spaces and probability laws in the arguments
below retain the tags as part of their codes.
\end{definition}

For a singleton image $\{s\}$ the left sides of
Equations~\eqref{eq:path-law-constraint} and
\eqref{eq:path-level-constraint} are bounded by $m_\xi(s)$ and $a_i(s)$,
respectively.  Countable subadditivity therefore shows that all the
constraints hold whenever the residual image is countable.  In
particular, ordinary finite-length paths in $G$ are allowed.  A simple
arc of finite domain-part length in an entire permitted level union
is also allowed: give it the one tag $(0,1)$, leaving at most its two
endpoints in the residual image.

\subsection{The compatible barrier theorem}

With the endpoint data and allowed path class specified, the construction
has the following quantitative form.

\begin{theorem}[Compatible barrier lemma]\label{thm:barrier}
Use any fixed path system of Definition~\ref{def:path-system} and
any endpoint data of Definition~\ref{def:endpoint}.  Suppose
$g\in L^2(G)$ is nonnegative Borel and bounded below by a positive
constant, and $m\in\ell^2(B)$ is nonnegative.  Assume that every
allowed path from $E$ to $b$ which first reaches $b$ at time $1$
and obeys the gates, if present, satisfies
\[
 \mathcal L^-_{g,m}(\gamma)\ge1.
\]
Then for every $\eps>0$ there is a barrier $u$, and an allowed
penalty region $O$ when gates are present, such that
\begin{equation}\label{eq:barrier-bound}
 \|\dd u\|_{2,G\setminus O}+\|a_u\|_{\ell^2(B)}
 \le C_{\mathrm{bar}}
       \bigl(\|g\|_2+\|m\|_{\ell^2(B)}\bigr)+\eps,
 \qquad C_{\mathrm{bar}}=8.
\end{equation}
One may take $0\le u\le1$, $a_u=4m$, and enlarge the flags only
by $\{m>0\}$.  The numerical constant is independent of the prior
list, earlier laws, atomless measure, null tuple exclusions, endpoint
data, and positive lower bound on $g$.
\end{theorem}

\subsection{Cutting, joining, and coding paths}

We prove the path properties needed for the theorem before constructing
the coordinate.  Cutting and joining must preserve the residual-contact
bounds; Borel coding will ensure measurability of the cost infimum.

\begin{lemma}[Restrictions and countable concatenations]
\label{lem:path-stability}
Allowed paths are stable under reversal, affine reparametrization,
and restriction to a closed parameter subinterval.  More generally,
suppose a continuous path is made from countably many restrictions of
allowed paths, possibly reversed, on disjoint open parameter intervals.
Assume that the complement of those intervals contains only countably
many times and that the resulting path has finite domain-part length.
Then the path can be tagged so as to be allowed.
\end{lemma}

\begin{proof}
For each of the two incidence inequalities, its left side is a
countably subadditive set functional of the residual image, continuous
under increasing unions.  Its value on $\{s\}$ is bounded by the
corresponding weight at $s$.  Consequently the inequalities for rational
closed intervals extend first to arbitrary open time intervals by an
increasing exhaustion.  They then extend to arbitrary closed intervals
by treating the two endpoints separately.

Restrict every old tag to the interior of the retained piece.  Keep
the restriction if the path is nonconstant there, and discard it
otherwise.  The residual image newly created by the discarded tags is
countable: each discarded restriction has one constant image point,
and there are only countably many tags.  This proves the restriction
claim, because countable image additions satisfy the incidence bounds
by singleton estimates and have zero $\lambda$ measure.  Reversal and
affine reparametrization are treated by pulling back the time intervals.

For the countable concatenation, keep these restricted tags in the
piece interiors.  The residual times fall into three groups: retained
old residual times, times exposed by discarded constant tags, and
junction or limit times.  Fix a test interval $I$ and a finite deletion
set $P$.  For each open piece interval $J$, let $R_J$ be the pullback of
the old residual times in $J\cap\operatorname{int}I$.  These times
are still residual, because every new tag is a restricted old tag.
Set
\[
 A_J=(Z_{\mathrm{new}}\cap J\cap\operatorname{int}I)\setminus R_J.
\]
The image of $A_J$ is countable: every such time belongs to a discarded
constant-tag restriction.  The remaining residual times $T$ consist
of junction or limit times, together with at most the two endpoints
of $I$, and hence are countable.  The sets $R_J$, $A_J$, and $T$ are
pairwise disjoint as sets of times.  The old bounds on open intervals
apply to each $R_J$.  If $A=T\cup\bigcup_J A_J$, its countable image
satisfies, for either incidence functional $\mathcal C$ and its
corresponding weight $w$,
\[
 \mathcal C(\gamma(A)\setminus P)
 \le\sum_{s\in\gamma(A)\setminus P}w(s)
 \le\sum_{\substack{r\in A\\\gamma(r)\notin P}}w(\gamma(r)).
\]
Apply the union bound on image sets and sum the old bounds with this
last one.  All charged times are disjoint, so their sum is bounded by
the required right side for the new path, even when their images
overlap.  The same decomposition proves the atomless condition.
Every retained tag has its old tuple and hence still avoids the
prescribed null set.

This proof explains the countability assumption on the extra times:
countably many open piece intervals alone could leave a Cantor set of
uncontrolled times.  In the collar concatenations below the intervals
accumulate only at explicitly specified dust endpoints, so the stated
assumption holds.
\end{proof}

To optimize over paths and place probability laws on their codes, we
also need the path constraints and costs to be Borel.

\begin{lemma}[Borel coding of paths and costs]\label{lem:path-coding}
For fixed path-system data, the allowed tagged codes form a standard
Borel space.  On this space the domain-part length and the arclength
occupation measure are Borel.  For every fixed nonnegative Borel
$g\in L^2(G)$ and fixed nonnegative $m\in\ell^2(B)$, all three costs
$\mathcal L_{g,m}$, $\mathcal L^-_{g,m}$, and
$\mathcal L^\circ_{g,m}$ are Borel, with values in $[0,\infty]$.
Endpoint restrictions, first arrival at a prescribed dust point at
time $1$, and hit or avoidance conditions for closed or relatively
open subsets of prescribed compact collar curves are Borel.
\end{lemma}

\begin{proof}
\emph{Tags and residual times.}
Start with $C([0,1],X)$ and the countable product of the slot spaces.
A slot is either inactive, or belongs to a countable disjoint union
of spaces consisting of a type, a level tuple, and endpoints
$0\le a<b\le1$.  These are standard Borel spaces.  Give the inactive
and active alternatives separate components.  Image containment for a
tag can be tested at rational times in $(a,b)$, using the closed level
graphs.  Nonconstancy can be tested by two such rational times with
different images.  Both requirements, as well as the tuple exclusion,
are Borel.  In this ambient coding the residual-time graph is closed
relative to the code space times $[0,1]$, and in particular has compact
time sections.

\emph{Incidence and visit counts.}
For Equation~\eqref{eq:path-law-constraint}, introduce the sampled path
$\eta$ and two time variables $r,u$.  The incidence in question is
equivalent to
\[
 r\in Z_\Gamma\cap I,\qquad
 \gamma(r)=\eta(u),\qquad \gamma(r)\notin P,
\]
for some $(r,u)\in[0,1]^2$.  Here $P$ is finite.  The last condition
is a countable union of closed conditions
$\dist_X(\gamma(r),P)\ge1/n$; for $P=\varnothing$ no condition is
needed.  Thus the witness sections are sigma-compact, and the
Arsenin--Kunugui projection theorem for sigma-compact sections applies
\cite[Theorem~18.18]{Kechris1995}.
For Equation~\eqref{eq:path-level-constraint}, use instead the closed
incidence relation $\gamma(r)\in H_i(v)$.  The same argument applies.
Integration against the fixed law $\xi$ or Haar measure shows that
the two left sides are Borel functions of the tagged code.  The test
$s\in\gamma(Z_\Gamma)$ is likewise Borel in $(\Gamma,s)$; integration
against $\lambda$ gives the Borel atomless condition.

For a fixed point $s$, at least $k$ visits in a specified closed
residual time set can be tested with $k$ time variables.  Distinctness
of those times is the union, over $n$, of the closed conditions
$|r_j-r_l|\ge1/n$ for $j\ne l$.  Projection again gives a Borel
visit count $N_s\in\{0,1,\ldots,\infty\}$.  Since every counting
$\ell^2$ weight has countable positive support,
\[
 \sum_{r\in\gamma^{-1}(B)}m(\gamma(r))
   =\sum_{s:m(s)>0}m(s)N_s(\gamma)
\]
is Borel.  The same argument treats the right sides of the two
constraints.  Omitting one or both endpoint times is handled by
exhausting the resulting open or half-open parameter intervals.

\emph{Physical arclength.}
For a rational closed interval $J$, the condition $\gamma(J)\subset G$
is open in path space.  On this set the physical restriction depends
continuously on the path, locally uniformly in a compact subset of $G$.
Its length is the supremum of the partition sums, and rational
partitions suffice by continuity.  Hence the length is Borel.  On
finite-length restrictions, the integral of a bounded nonnegative
continuous function $f$ against arclength is the limit of sums
\[
 \sum_j f(\gamma(t_j))
       \length_{\mathrm{sph}}(\gamma|_{[t_j,t_{j+1}]}),
\]
over increasingly fine equal subdivisions.  Uniform continuity of
$f\circ\gamma$ and finiteness of total length justify this limit.
The full domain-part integral is the supremum of sums over finite
disjoint collections of rational closed intervals contained in the
open domain times.  Such collections exhaust all those times, and
arclength has no atoms.  This proves the Borel length assertion.
A countable determining family of continuous functions on $G$ also
shows that $\gamma\mapsto\nu_\gamma$ is Borel into the space of
finite Borel measures on $G$.  Consequently
$\int g\,\dd\nu_\gamma$ is Borel for every fixed nonnegative Borel
$g$, by the usual monotone-class argument, including extended values.

The finite-length condition and all the constraints are therefore
Borel, proving the assertion about the tagged space.  Endpoint
conditions are immediate.  Avoidance of a fixed point on $[0,1)$
can be tested on the compact intervals $[0,1-1/n]$.  Hit conditions
for closed subsets of a compact collar have compact time sections;
relatively open collar subsets are sigma-compact and give the same
Borel conclusion.  These observations prove the final assertion.
\end{proof}

\begin{remark}
The preceding lemma concerns tagged codes.  The set of untagged paths
that \emph{admit} tags is an analytic projection and need not be Borel.
Likewise, cost measurability is asserted for each fixed metric $(g,m)$;
no Borel parametrization of all counting $\ell^2(B)$ weights is being
asserted.  All probability measures on analytic sets may be completed,
as in Section~\ref{sec:setup}.
\end{remark}

\subsection{Collar concatenation}

The following lemma converts bounded-cost paths approaching dust
endpoints into paths that reach them.  It is the compactness substitute
used twice in the barrier construction: assembling actual path pieces
lets us retain the path restrictions through the stability lemma.

\begin{lemma}[Collar concatenation]\label{lem:collar-splicing}
Fix a path system and a metric $(g,m)$ with $g\ge c_0>0$ on $G$.
Let allowed paths $\gamma_n$, whose endpoints are in $G$, have costs
converging to a finite number $A$.
\begin{enumerate}[label=\textup{(\roman*)}]
\item If their initial and final endpoints converge to distinct points
$p,q\in B$, then for every $\eta>0$ there is an allowed path $\gamma$
from $p$ to $q$ satisfying
\[
 \mathcal L^\circ_{g,m}(\gamma)\le A+\eta.
\]
\item If their initial endpoints belong to a set $F\subset G$ and
their final endpoints converge to $p\in B\setminus\overline F$,
where closure is in $X$, there is an allowed path from a point of $F$
to $p$ with $\mathcal L^-_{g,m}\le A+\eta$.
\end{enumerate}
The paths produced consist of countably many restricted original
pieces and joining arcs on smooth collar curves in $G$.  Every dust
point avoided by the original pieces remains avoided away from the
newly adjoined endpoint times.
Any prescribed hit restrictions are also preserved when the joining
collars are disjoint from their hit curves.
\end{lemma}

\begin{proof}
We take limits of collar hit positions and fragment costs, then join
actual fragments with summable extra cost.  Shrinking disks will give
continuity at the new endpoint times.

\emph{Choice of finite-cost collars.}
About $p$ and $q$ choose strictly nested dust disks $V_j$ and $W_j$
with closures shrinking to their respective points and with disjoint
outer closures.  Their boundary curves $A_j=\partial V_j$ and
$B_j=\partial W_j$ can be chosen to satisfy
\[
 \int_{A_j}g\,\dd s<\infty,
 \qquad \int_{B_j}g\,\dd s<\infty.
\]
Indeed each smooth boundary has a compact smooth annular collar in
$G$, and Fubini in its smooth parallel coordinates gives finite
integrals for almost every parallel curve.  Arbitrarily small
perturbations preserve all the required nesting and separation.
For part \textup{(ii)}, choose only the collars about $p$, with the
outer closed disk missing $\overline F$.  When joining collars
disjoint from additional hit curves are available, choose them.
For the nested gate curves used below, the intervening annuli
provide these choices.

\emph{Cuts and limiting budgets.}
Consider part \textup{(i)}.  After passage to a subsequence the two
endpoints of $\gamma_n$ lie in $V_n$ and $W_n$.  For each fixed
$j\le n$, let $\alpha_{n,j}$ be its first hit of $A_j$ and let
$\beta_{n,j}$ be its last hit of $B_j$.  The cuts are ordered as
\[
 \alpha_{n,n}\le\cdots\le\alpha_{n,1}
       <\beta_{n,1}\le\cdots\le\beta_{n,n}.
\]
The strict middle inequality follows from disjointness of the outer
closed disks.  The initial fragment between $\alpha_{n,j+1}$ and
$\alpha_{n,j}$ lies in $\overline V_j$, and the final fragment
between $\beta_{n,j}$ and $\beta_{n,j+1}$ lies in $\overline W_j$.
Retain also the middle fragment between $\alpha_{n,1}$ and
$\beta_{n,1}$.

Every cut point belongs to $G$.  Thus these fragments have disjoint
interior time sets, no dust toll is duplicated at a cut, and arclength
adds without endpoint atoms.  A diagonal subsequence makes all the
hit positions and all the fragment costs converge.  Denote the
limiting middle cost by $c_*$ and the two families of limiting
fragment costs by $c_j^-$ and $c_j^+$.  Every finite partial sum is
bounded by $A$, so
\begin{equation}\label{eq:collar-cost-budget}
 c_*+\sum_{j\ge1}(c_j^-+c_j^+)\le A.
\end{equation}

\emph{Joining actual fragments.}
On each fixed collar, two points sufficiently close to the same
limiting hit position can be joined by a collar arc whose $g$ integral
is arbitrarily small.  This is absolute continuity of the integral
of $g$ along that finite-integral smooth circle.  Assign summable
positive error allowances, with total less than $\eta$, both to
fragment costs and to all joining arcs.  At each collar choose a
neighborhood of its limiting hit position small enough for the
assigned joining allowance.  Select the middle fragment and every
initial and final fragment from sufficiently far along its convergent
subsequence so that its cost is at most its limit plus its allowance,
and both endpoints are in the chosen collar neighborhoods.  Adjacent
selected fragments can then be joined with the assigned small costs.

Parametrize the resulting initial pieces on successive intervals
accumulating at $0$, the middle portion on a compact interior interval,
and the final pieces on successive intervals accumulating at $1$.
The initial tails lie in shrinking $\overline V_j$ and the final tails
in shrinking $\overline W_j$.  Adding endpoint values $p,q$ therefore
gives a continuous path in $X$.  Equation~\eqref{eq:collar-cost-budget}
and the assigned errors bound its cost with the two endpoint tolls
omitted by $A+\eta$.  Since $g\ge c_0$, the sum of all domain-part
lengths is finite.  There are only countably many additional junction
times.  Lemma~\ref{lem:path-stability} makes the new path allowed.

For part \textup{(ii)}, use last hits of the collars about $p$.  Keep
one actual prefix, beginning at its original point of $F$, followed
by the successive last-hit fragments.  The same diagonal argument,
budget inequality, and small collar arcs construct the required
one-ended path.  Finally, the construction uses only restrictions
of original paths and arcs on the selected collars.  This proves
the asserted preservation of avoidance and hit restrictions.
\end{proof}

\subsection{Removing a newly introduced level of small jumps}

Jumps inside small dust disks will control the eventual barrier's
oscillation near the dust.  We must show that adding these jumps
preserves a positive cost for travel from an anchor to a collar.

A \emph{jump level} consists of a finite pairwise disjoint cover of
$B$ by dust disks whose closures miss the anchor $E$ and a specified
collar $C$.  Each disk has a positive toll, independent of the two
endpoints of a jump.  A jump joins any two domain points inside that
disk.  A finite chain alternates allowed true paths with jumps; all
true-path endpoints and all jump endpoints belong to $G$.
Its cost is the sum of true costs and jump tolls.  A constant true
path in $G$ is allowed, so consecutive jumps cause no difficulty.

\begin{lemma}[Removal of one shrinking jump level]
\label{lem:jump-removal}
Fix a metric $(g,m)$ with $g\ge c_0>0$, a compact anchor $E\subset G$, and
a compact collar curve $C\subset G$.  Suppose that, with finitely
many fixed jump levels, every chain from $E$ to $C$ has cost at least
$\beta>0$.  Fix $\rho,\tau>0$ and $0<\Delta<\beta$.
For all sufficiently small meshes, adding any new jump level whose
disk $V$ has toll
\begin{equation}\label{eq:neighborhood-jump-cost}
 c(V)=\max\left\{\rho,
       \sup_{\substack{s\in B\\\dist_X(s,V)<\tau}}m(s)\right\}
\end{equation}
preserves the lower bound $\beta-\Delta$ for chains from $E$ to $C$.
Here each new disk meets $B$ and its closure misses $E$ and $C$.
\end{lemma}

\begin{proof}
In the contradiction argument, shrinking new jumps collapse to dust
junctions.  Collar concatenation replaces the adjoining path pieces,
and the new-jump tolls pay for the junction visits.

\emph{Fixing the jump pattern and its limiting costs.}
Suppose instead that covers of mesh tending to zero admit chains of
cost less than $\beta-\Delta$.  All existing jump tolls and $\rho$
are bounded below by a common positive number.  The number of jumps
in these chains is therefore uniformly bounded.  After combining
successive true paths and passing to a subsequence, fix the number
of jumps, the pattern of old and new jumps, and the disk label at
every old jump.  The two endpoints of each new jump converge, after
a further subsequence, to the same point $p_i\in B$: its disk meets
$B$ and has diameter tending to zero.  All true-piece costs and all
new-jump costs may also be assumed convergent.  Since $\tau$ is fixed,
the limiting cost of the new jump at $p_i$ is at least $m(p_i)$.

\emph{Replacing the blocks between old jumps.}
Break each chain into blocks separated by old jumps.  The anchors
of a block are $E$, $C$, or the open disk of the adjacent old jump.
We will replace the block by an allowed true path between its two
anchors, with cost at most the sum of its limiting costs plus
arbitrarily small error.  Different blocks need not choose the same
domain point inside an old jump disk: that old jump can still join
their chosen endpoints at its fixed toll.

Suppose first that a block has new-jump limits
$p_1,\ldots,p_d$.  Collapse each run of equal consecutive limits
to one junction and discard the intervening true pieces.  This only
decreases the available budget; one of the jump budgets in the run
still pays at least the required single toll.  Between successive
distinct junctions apply Lemma~\ref{lem:collar-splicing}\textup{(i)}
to the corresponding true pieces.  Their endpoints lie in $G$ and
converge to those junctions.  This produces an allowed path between
them at its limiting true-piece cost plus an arbitrary small error,
with the two junction tolls temporarily omitted.

We must also attach the replacement to its two outer anchors, with
endpoints in the domain.
At an outer anchor $A$, if its adjacent junction $p$ is outside
$\overline A$, use the one-ended version of
Lemma~\ref{lem:collar-splicing}, reversing paths if necessary.  If
$p\in\overline A$, the anchor cannot be $E$ or $C$, since these
are compact subsets of $G$.  It is therefore an old jump disk.
Its boundary lies in $G$, so $p\in B\cap\overline A$ actually
belongs to its interior.  In this case discard the outer true
piece and truncate the next constructed path at a domain point
sufficiently close to $p$ to lie in $A$.  Such points exist
arbitrarily close to $p$ in the collar construction.  More generally,
a nonconstant arc cannot have an initial interval contained in the
totally disconnected set $B$.  The truncation is allowed by
Lemma~\ref{lem:path-stability} and cannot increase the cost.

If a single remaining junction lies in both anchors, use instead a
common domain point of their open intersection and the constant path.
Such a point exists because $G$ is dense in $X$.  If precisely one
anchor contains the sole remaining junction, construct the one-ended
path from the other anchor to that junction and truncate its end on
a sufficiently deep domain collar inside the first anchor.  The
truncation adds no dust toll, and reversal gives the required order.
If the block has no new junctions, simply retain an actual true path
with cost sufficiently close to its limit.

\emph{Accounting for the dust junctions and reconnecting the blocks.}
Concatenate the retained constructed paths at each remaining dust
junction using a single shared time, without adding a pause there.
Exactly one previously omitted toll is then incurred at that time.
It is covered by the limiting new-jump cost, which is at least
$m(p_i)$.  Any additional internal visits to that point were already
included in the true-piece costs.  Infinite collar tails add only
their specified endpoint times, and the construction uses countably
many pieces in all.  Its domain-part length is finite because
$g\ge c_0$.  Thus Lemma~\ref{lem:path-stability} supplies the
required allowed true path replacing the block.

There are finitely many blocks.  Choose their errors to have total
less than $\Delta/2$ and reconnect them with the unchanged old jumps.
The resulting old-level chain has cost at most
$\beta-\Delta+\Delta/2<\beta$, a contradiction.
\end{proof}

\subsection{Proof of the compatible barrier theorem}

\begin{proof}[Proof of Theorem~\ref{thm:barrier}]
Steps 1 and 2 reduce the endpoint hypothesis to a positive cost for
reaching a fixed collar, with the relevant finitely many gates absorbed
into a density on a penalty region.  Steps 3--5 construct a separating
function from chain costs, bound its topological bands, and smooth it
with controlled energy while retaining those band bounds.  Step 6
proves compatibility with the prior coordinates by applying the path-cost estimate to
tagged arcs in their level unions.

\emph{Step 1: a finite collar already has positive cost.}
We claim that for some $k$, every allowed path from $E$ to $C_k$
obeying the gates through index $k$ has full cost at least $1/2$.
Otherwise choose violations with indices $k_n\to\infty$ and cut each at its first
hit of $C_{k_n}$.  The resulting paths do not hit $b$ or any deeper
prescribed collar.  They consequently obey all the gates, including
those with index greater than $k_n$, since those have no hits.
Their terminal points converge to $b$.  After taking a subsequence
their costs converge to a number at most $1/2$.

Choose strictly intermediate smooth collars between the prescribed
ones, with finite $g$ integrals.  They avoid every prescribed hit
curve.  Apply the one-ended construction in
Lemma~\ref{lem:collar-splicing} with error less than $1/4$.
It gives an allowed path from $E$ to $b$ with final toll omitted
and cost less than $1$.  The retained original pieces avoid $b$,
the joining arcs lie in $G$, and all gates are preserved.  Thus its
first visit to $b$ is at the final endpoint, contrary to hypothesis.

\emph{Step 2: replace finitely many gates by a penalty density.}
In the ungated case set $O=\varnothing$ and $g_1=g$.
With gates, form the compact set
\[
 F=\bigcup_{j=1}^k\bigl(C_j\setminus f_0^{-1}(W_2)\bigr)
       \subset G.
\]
It is disjoint from $E$, and its image misses $W_1$.  Compact
separation gives a relatively compact open set $O\subset G$ containing
$F$ whose closure misses $E$ and whose image under $f_0$ misses $W_1$.
Choose a bounded nonnegative density $h$ vanishing outside $O$ such that
every path from $E$ to $F$ has $h$-length at least $1$.  For example,
if $F\ne\varnothing$, put
\[
 \delta=\dist_{\mathrm{sph}}(F,\Chat\setminus O)>0,
 \qquad h=\delta^{-1}\ind_O.
\]
The last passage from outside $O$ to $F$ is entirely in $G$ and has
physical length at least $\delta$.  If $F=\varnothing$, take $h=0$
and $O=\varnothing$.  Set $g_1=g+h$.  This density is in $L^2(G)$,
remains bounded below by a positive constant, and agrees with $g$
off $O$.

An allowed path from $E$ to $C_k$ which fails a gate through index
$k$ reaches $F$ and has $g_1$-cost at least $1$.  A path which obeys
those gates has the $1/2$ lower bound from Step 1.  Cutting first
at $C_k$ if necessary shows that every allowed path from $E$ to
$C_k$, without any gate restriction, now has full $(g_1,m)$-cost
at least $1/2$.

\emph{Step 3: chain distance with progressively smaller jumps.}
Taking an infimum of path costs is the classical modulus-to-capacity
construction; see \cite[Proposition~2.17]{HeinonenKoskela1998}.
Here shrinking jumps will control the bands at every dust point.
Choose positive losses $\Delta_j$ with $\sum_j\Delta_j<1/8$.
Starting with the true paths alone, add jump levels one at a time.
At level $j$ take $\rho_j,\tau_j\downarrow0$ and a dust-disk cover
of mesh tending to zero, with closures missing $E$ and $C_k$.
Lemma~\ref{lem:dust-disks} permits these covers, and
Lemma~\ref{lem:jump-removal} permits choosing each new mesh so small
that the lower bound decreases by at most $\Delta_j$.  Thus all
finite chains using any of the levels have cost at least
\[
 \beta_*=\frac12-\sum_j\Delta_j>\frac14
\]
from $E$ to $C_k$.  Each such chain uses only finitely many levels,
so this assertion follows from the corresponding finite-stage bound.

For $x\in G$ let $D(x)$ be the infimum of the costs of these finite
chains from $E$ to $x$, allowing the value $\infty$, and define
\begin{equation}\label{eq:barrier-chain-distance}
 u^0(x)=\min\{1,4D(x)\}.
\end{equation}
The empty chain shows $u^0=0$ on $E$.  A jump disk cannot cross
$C_k$, because it is connected and its closure misses that curve.
In a chain from $E$ to a point of $U_k\cap G$, some true-path piece
must therefore cross $C_k$.  Truncating that chain at the crossing
gives cost at least $\beta_*$.  Consequently
\begin{equation}\label{eq:barrier-constant-regions}
 u^0=0\text{ on }E,
 \qquad u^0=1\text{ on }U_k\cap G.
\end{equation}
Reversal and concatenation give, for every allowed true path
from $x\in G$ to $y\in G$,
\begin{equation}\label{eq:barrier-path-inequality}
 |u^0(x)-u^0(y)|\le4\mathcal L_{g_1,m}(\gamma).
\end{equation}
This remains valid when one or both distances are infinite: a
finite-cost connecting path makes finiteness of the two distances
equivalent, and otherwise the inequality is vacuous or both truncated
values are $1$.  Likewise any two domain points of a level-$j$ jump
disk $V$ have value difference at most $4c_j(V)$.

\emph{Step 4: topological bands.}
For each $s\in B$, the level-$j$ disk $V_j(s)$ containing $s$ contains
all domain points sufficiently close to $s$.  Its oscillation bound
therefore bounds the width of the topological band at $s$ by
$4c_j(V_j(s))$.  We claim
\[
 \limsup_{j\to\infty}c_j(V_j(s))\le m(s).
\]
Indeed, for every $\eta>0$ the set
$\{r\in B:m(r)\ge m(s)+\eta\}$ is finite and does not contain $s$.
As the mesh and $\tau_j$ tend to zero, its points eventually lie
outside the $\tau_j$-neighborhood of $V_j(s)$.  Also $\rho_j\to0$.
Equation~\eqref{eq:neighborhood-jump-cost} proves the claim.
Hence every band has width at most $4m(s)$.  In particular $u^0$
has a unique trace at every point of
\[
 B\setminus\Sigma',\qquad
 \Sigma'=\Sigma\cup\{m>0\}.
\]
This is a countable enlargement of the flags.

\emph{Step 5: Sobolev estimates and band-preserving smoothing.}
The separating values and pointwise band bounds for $u^0$ are now
established.  It remains to obtain smoothness and the energy estimate.

\emph{Measurability and weak derivatives.}
The finite-chain codes retain all true-path tags and are Borel by
Lemma~\ref{lem:path-coding}:
matching endpoints, anchor membership, and membership in specified
jump disks are Borel conditions, and their costs are Borel.
Projection over such codes shows that $\{x:D(x)<a\}$ is analytic
for every $a$.  Thus $u^0$ is measurable for completed area measure.

In a physical conformal chart, write $\dd s=\Lambda(z)|\dd z|$.
Ordinary axis segments in the chart are allowed and contain no dust
tolls.  Equation~\eqref{eq:barrier-path-inequality} gives
\[
 |u^0(z_2)-u^0(z_1)|
   \le4\int_{[z_1,z_2]}\Lambda g_1\,|\dd z|
\]
on every such segment.  Fubini makes the line density integrable
on almost every horizontal and vertical line in each compact chart
rectangle.  On each of these lines the displayed inequalities for
all subsegments imply absolute continuity.  Applying the usual
slicing characterization of weak derivatives to an almost everywhere
equal Borel representative gives
\[
 |\partial_x u^0|,|\partial_y u^0|\le4\Lambda g_1,
 \qquad |\dd u^0|_{\mathrm{sph}}\le4\sqrt2\,g_1
 \quad\text{almost everywhere}.
\]
Since $0\le u^0\le1$ and spherical area is finite, this is a global
$W^{1,2}(G)$ function.

\emph{Smoothing with shrinking averaging neighborhoods.}
Choose a locally finite smooth partition of unity on relatively compact
coordinate patches, and slightly larger patches in which to take
positive mollified averages.  The patches
can be chosen so that every averaging sample contributing at a
point $x$ has $X$-distance from $x$ tending to zero as $x$ approaches
$B$.  For example, first require their $X$ diameters to be at most
a fixed small multiple of their distance to $B$, and then shrink
the coordinate averaging radii further.  On each patch approximate
$u^0$ by a positive convolution average.  Choose its value and
derivative errors in $L^2$ so small that, including the derivatives
of the partition functions, the sum of all errors is less than
$\eps$.  Local Sobolev approximation permits this on each of the
countably many relatively compact patches.  The resulting smooth
function $u$ satisfies
\[
 0\le u\le1,
 \qquad \|\dd u-\dd u^0\|_{2,G}<\eps.
\]
For precision, if $\psi_j$ is the partition and $v_j$ the local
average, estimate
\[
 \dd u-\dd u^0
   =\sum_j\psi_j(\dd v_j-\dd u^0)
       +\sum_j(v_j-u^0)\dd\psi_j
\]
by the triangle inequality in $L^2$, assigning summable tolerances
to both terms.

\emph{Preservation of bands and constants.}
To verify the band assertion pointwise at every dust point, let
$\ell_s,r_s$ be the actual lower and upper topological limits of
$u^0$ at $s\in B$.  For every $\eta>0$ all actual values of $u^0$
in some punctured domain neighborhood of $s$ lie in
$[\ell_s-\eta,r_s+\eta]$.  Changing to an almost everywhere equal
Borel representative preserves these bounds almost everywhere in
that neighborhood.  If $x\to s$ and a sample $y$ contributes to
the average at $x$, the patch construction gives, for a fixed
constant $c$,
\[
 d_X(y,s)\le d_X(x,s)+c\,\dist_X(x,B)\longrightarrow0,
\]
uniformly over the contributing patches and samples.  Eventually
all those averages, and then their positive partition sum, lie in
the indicated interval.  Letting $\eta\downarrow0$ proves that
all limiting values lie in the original band at $s$, without an
exceptional dust set.
Thus the smoothed bands are contained in those of $u^0$, and every
trace at a new free point is unchanged.  Equation~\eqref{eq:barrier-constant-regions}
also permits exact preservation of the required constants.  Choose
$l>k$ and arrange that every sampling patch meeting $E_*$ lies in
$\operatorname{int}E$, while every sampling patch meeting $C_l$
lies in $U_k\cap G$.  The two compact sets are disjoint and have
the requisite open neighborhoods.  Their averages are exactly $0$
and $1$, respectively.

\emph{The controlled size.}
Take $a_u=4m$.  Since $g_1=g$ off $O$,
\[
 \|\dd u\|_{2,G\setminus O}+\|a_u\|_{\ell^2(B)}
 \le4\sqrt2\,\|g\|_2+4\|m\|_{\ell^2(B)}+\eps
 \le8\bigl(\|g\|_2+\|m\|_{\ell^2(B)}\bigr)+\eps.
\]
Total energy is finite as well, because $g_1\in L^2(G)$.

\emph{Step 6: compatibility with every prior finite level union.}
Only \textup{(P)} for the appended coordinate remains.  We will show
that its free trace values on each nondegenerate component of a generic
prior finite level union form a null set.

Fix a type $\boldsymbol\iota$ and choose a tuple outside
$N_{\boldsymbol\iota}$ and the null sets in
Lemma~\ref{lem:level-continua} for the fixed weights $g_1,m$.  Put
$H=H_{\boldsymbol\iota}(\mathbf t)$.  By
Lemma~\ref{lem:level-continua} and its weighted coarea estimate,
$H$ has finite physical length in $G$, its nondegenerate components
are locally connected and countable in number, and
\begin{equation}\label{eq:barrier-level-measure}
 \nu_H
   =\ind_{H\cap G}g_1\,\dd\HH^1_{\mathrm{sph}}
          +\sum_{s\in H\cap B}m(s)\delta_s
\end{equation}
is a finite Borel measure.  The integrability behind this last
assertion is, explicitly, bounded by the finite sum of
\[
 \int_{G\cap M_i}g_1|\dd u_i|\,\dd a
       +\sum_{s\in B}m(s)a_i(s)
\]
over the participating labels.  Cauchy--Schwarz bounds each term.
This also treats extended nonnegative values of $g_1$, by truncation.

Let $\alpha\subset H$ be a simple arc joining two new free points
$s,t\in B\setminus\Sigma'$.  Its domain-part arclength is its
physical image length and is finite.  Give it the tag corresponding
to the fixed permitted tuple.  Subarcs between domain points are
allowed, so Equation~\eqref{eq:barrier-path-inequality} bounds their
value differences by four times their $\nu_H$ measure.  Domain
points occur arbitrarily close to both ends of $\alpha$, since $B$
is totally disconnected.  Passing to the free trace limits gives
\[
 |u^0(s)-u^0(t)|\le4\nu_H(\alpha).
\]
The same trace inequality holds for $u$, whose free traces were
unchanged by smoothing.  The measure $\nu_H$ gives zero mass to
the new free dust points.  Apply Lemma~\ref{lem:trace-null} in
each nondegenerate component of $H$.  The set of its free trace
values has Lebesgue outer measure zero, hence zero image measure
modulo $1$.  There are only countably many such components.

For the fixed generic prior tuple, almost every new level consequently
contains no free point in a nondegenerate component of $H$.
The universal measurability of the level-continuum incidence from
Section~\ref{sec:setup} permits slicing in the tuple variables.
This is exactly \textup{(P)} for the appended coordinate.  Older
instances of \textup{(P)} persist because enlarging the flags only
removes free points.  The barrier and all its asserted bounds are
proved.
\end{proof}
